\documentclass[aip,jcp,reprint,noshowkeys,superscriptaddress]{revtex4-1}
\usepackage{graphicx,dcolumn,bm,xcolor,microtype,multirow,amscd,amsmath,amssymb,amsfonts,physics,longtable,wrapfig}
\usepackage[version=4]{mhchem}

\usepackage[utf8]{inputenc}
\usepackage[T1]{fontenc}
\usepackage{txfonts}

\usepackage[
    colorlinks,
    linkcolor={red!50!black},
    citecolor={red!70!black},
    urlcolor={red!80!black}
	]{hyperref}
\urlstyle{same}

\newcommand{\ie}{\textit{i.e.}}
\newcommand{\eg}{\textit{e.g.}}
\newcommand{\alert}[1]{\textcolor{red}{#1}}
\usepackage[normalem]{ulem}
\newcommand{\titou}[1]{\textcolor{red}{#1}}
\newcommand{\trashPFL}[1]{\textcolor{\red}{\sout{#1}}}
\newcommand{\PFL}[1]{\titou{(\underline{\bf PFL}: #1)}}

\newcommand{\cre}[1]{\Hat{a}_{#1}^{\dag}}
\newcommand{\ani}[1]{\Hat{a}_{#1}}
\newcommand{\ca}[2]{\Hat{a}^{#1}_{#2}}

\newcommand{\mc}{\multicolumn}
\newcommand{\fnm}{\footnotemark}
\newcommand{\fnt}{\footnotetext}
\newcommand{\tabc}[1]{\multicolumn{1}{c}{#1}}
\newcommand{\QP}{\textsc{quantum package}}
\newcommand{\T}[1]{#1^{\intercal}}

% coordinates
\newcommand{\br}{\boldsymbol{r}}
\newcommand{\bx}{\boldsymbol{x}}
\newcommand{\bI}{\boldsymbol{1}}
\newcommand{\cK}{\mathcal{K}}

% orbital energies
\newcommand{\eps}{\epsilon}
\newcommand{\irr}{\text{irr}}

% shortcuts for greek letters
\newcommand{\si}{\sigma}
\newcommand{\la}{\lambda}
\newcommand{\ii}{\mathrm{i}}
\newcommand{\up}{\uparrow}
\newcommand{\dw}{\downarrow}
\newcommand{\upup}{\up\up}
\newcommand{\updw}{\up\dw}
\newcommand{\dwup}{\dw\up}
\newcommand{\dwdw}{\dw\downarrow}

\newcommand{\h}{\text{1h}}
\newcommand{\p}{\text{1p}}
\newcommand{\hp}{\text{1h+1p}}
\newcommand{\hhp}{\text{2h1p}}
\newcommand{\pph}{\text{2p1h}}
\newcommand{\ph}{\text{ph}}
\newcommand{\eh}{\text{eh}}
\newcommand{\he}{\text{he}}
\newcommand{\ee}{\text{ee}}
\newcommand{\teh}{\overline{\text{eh}}}
\newcommand{\pp}{\text{pp}}
\newcommand{\hh}{\text{hh}}
\newcommand{\xc}{\text{xc}}
\newcommand{\Hx}{\text{Hx}}

% addresses
\newcommand{\LCPQ}{Laboratoire de Chimie et Physique Quantiques (UMR 5626), Universit\'e de Toulouse, CNRS, UPS, France}

\begin{document}

\title{Multiresolvent Representation of Frequency-Dependent Vertices}

%\author{Pierre-Fran\c{c}ois \surname{Loos}}
%	\email{loos@irsamc.ups-tlse.fr}
%	\affiliation{\LCPQ}

\begin{abstract}
A general four-point vertex depends on three independent frequencies and requires a richer resolvent structure than those used in one-body Green's function, self-energy, and polarization-propagator theories.
We develop a channel-resolved multiresolvent construction of the parquet vertex through third order in the bare interaction, taking the electron-hole ($\eh$) and particle-particle ($\pp$) reducible vertices as primitive objects and generating the transverse electron-hole contribution by crossing symmetry.
For a Hartree--Fock reference, an exact factorization of the non-interacting pair propagator defines one-sided frequency projections and left and right coupling vertices.
At second order, the couplings are instantaneous and the frequency dependence is carried by central resolvents associated with neutral $N$-electron and $(N\pm2)$-electron propagation.
At third order, the projected couplings acquire $\hhp$ and $\pph$ poles, while the dynamic irreducible kernels activate additional orbital sectors: occupied-occupied and virtual-virtual propagation in the $\eh$ channel, and mixed occupied-virtual propagation in the $\pp$ channel.
The complete third-order result retains these contributions alongside the ordinary coupling and central-resolvent corrections.
Detailed derivations are given in two channel-specific appendices.
Second-order consistency checks recover the corresponding single-frequency GF2 kernels and, on applying the same reduction to the conventional dynamical $GW$ kernel at first order in $W$, the known dynamical $\eh$ and $\pp$ Bethe--Salpeter kernels.
\end{abstract}

\maketitle

%\tableofcontents

%%%%%%%%%%%%%%%%%%%%%%%%%
\section{Introduction}
\label{sec:introduction}
%%%%%%%%%%%%%%%%%%%%%%%%%

A general four-point quantity contains substantially more analytic
structure than a one-particle Green's function or a linear-response
propagator.
For a time-independent Hamiltonian, time-translation invariance removes
one overall time variable, leaving three independent time differences or,
equivalently, three independent frequencies.
The full two-particle vertex can therefore be written as
$F(\nu,\nu',\omega)$.\cite{StarkeKresse2012}

This three-frequency dependence makes a direct extension of conventional
spectral or algebraic-diagrammatic construction (ADC) non-trivial.
A one-particle self-energy can be represented by resolvents in the
$(N\pm1)$-electron sectors, while a polarization propagator is naturally
expressed through a neutral $N$-electron resolvent.
For a four-point vertex, several propagation sectors and several independent
frequencies coexist.
The objective is therefore not to introduce a single formal resolvent
carrying all three frequencies, but to organize the vertex into ordinary
resolvents associated with well-defined intermediate-state spaces.

Parquet theory provides the natural channel decomposition for this purpose.
We denote the electron-hole, transverse electron-hole, and particle-particle
channels by $\eh$, $\teh$, and $\pp$, respectively.
The full vertex is decomposed as\cite{DeDominicis_1964b,Bickers_2004,Eckhardt2023,KuglerVonDelft2018,MarieLoos2025}
\begin{equation}
	F
	=
	\Lambda
	+
	\Phi^\eh
	+
	\Phi^{\teh}
	+
	\Phi^\pp
	\label{eq:parquet_decomposition}
\end{equation}
where $\Lambda$ is fully two-particle irreducible.
The direct $\eh$ reducible vertex $\Phi^\eh$ is associated with neutral
$N$-electron propagation, whereas $\Phi^\pp$ is associated with the
$(N\pm2)$-electron sectors.
The transverse contribution $\Phi^{\teh}$ follows from $\Phi^\eh$ by
fermionic crossing.
Consequently, only two primitive multiresolvent constructions are required:
$\Phi^\eh$ and $\Phi^\pp$.

In each primitive channel, the ordinary channel-pair contribution is
organized as a left coupling vertex, a channel-specific central resolvent,
and a right coupling vertex.
The central resolvent carries the natural bosonic transfer of the channel,
while the coupling vertices retain the remaining fermionic-frequency dependence.
At third order, the frequency-dependent irreducible kernels also activate
orbital sectors outside the ordinary channel-pair spaces; these contributions
are retained explicitly.
This construction also provides a direct route to the four-point Hedin
kernel.
The vertex irreducible in the direct $\eh$ channel is
\begin{equation}
	\Gamma^\eh
	=
	\Lambda
	+
	\Phi^{\teh}
	+
	\Phi^\pp
	\label{eq:Gamma_eh_parquet}
\end{equation}
At the exact level, this is the same four-point kernel that enters Hedin's
Bethe--Salpeter equation (BSE) through the functional derivative
$\Xi=\fdv*{\Sigma}{G}$, once a common index convention is adopted 
.~\cite{Hedin1965,Onida2002,StarkeKresse2012,MarieAmmarLoos2024}
Thus, constructing $\Phi^\eh$ and $\Phi^\pp$ also yields a multiresolvent
representation of the direct $\eh$ Hedin kernel.

The present notes establish the primitive channel construction through
second and third order in the bare interaction.
The second-order result involves instantaneous couplings and non-interacting
central resolvents. At third order, corrections to both couplings and central
propagation coexist with additional contributions from activated orbital sectors.
The main text presents the complete structure, while Appendices~\ref{sec:eh_third}
and \ref{sec:pp_third} give the detailed frequency integrations.
Independent second-order GF2 and dynamical-$GW$ reductions are summarized
in Sec.~\ref{sec:channel_consistency}.

Throughout the perturbative discussion, we use a Hartree--Fock (HF) zeroth-order
reference together with a M{\o}ller--Plesset (MP) partition.
In the one-shot construction, $G=G_0$; accordingly, $G_0$ denotes the non-interacting one-particle propagator and
$L_0=G_0G_0$ the corresponding non-interacting two-particle propagator.
For compactness, all expressions are written for real-valued orbitals.
Left and right coupling vertices are kept distinct whenever this distinction
matters, and all causal prescriptions are denoted by $0^+$.

%%%%%%%%%%%%%%%%%%%%%%%%%
\section{Frequency conventions and channel organization}
\label{sec:frequencies}
%%%%%%%%%%%%%%%%%%%%%%%%%


Occupied spinorbitals are denoted by $i,j,k,l$, virtual spinorbitals by
$a,b,c,d$, and general spinorbitals by $p,q,r,s$.
Antisymmetric particle-pair and hole-pair states are represented by unique
orbital pairs, $a<b$ and $i<j$, respectively.
All matrix contractions in these auxiliary spaces use unique pairs, and the
coupling vertices contain no $1/\sqrt{2}$ normalization factors.
The two relative-frequency orientations of a unique pair are retained and
averaged when defining frequency projections.
Orbital pairs describing electron-hole propagation remain unrestricted.
Whenever an explicit sum represents an antisymmetric particle or hole pair,
its ordering restriction is displayed.

The direct $\eh$ convention is used throughout these notes.
The natural bosonic transfer frequencies of the three parquet channels are\cite{OrlandoRomanielloLoos2023,MarieLoos2025}
\begin{align}
	\Omega_\eh
	&=
	\omega
	\label{eq:Omega_eh}
	\\
	\Omega_{\teh}
	&=
	\nu'-\nu
	\label{eq:Omega_ehbar}
	\\
	\Omega_\pp
	&=
	\omega+\nu+\nu'
	\label{eq:Omega_pp}
\end{align}

For the transverse $\eh$ channel, exchanging the third and fourth external
legs gives the frequency map
\begin{equation}
	\left(
	\nu_{\teh},
	\nu'_{\teh},
	\Omega_{\teh}
	\right)
	=
	\left(
	\nu,
	\nu+\omega,
	\nu'-\nu
	\right)
	\label{eq:crossed_frequency_map}
\end{equation}
The natural $\pp$ variables are related to the direct $\eh$ convention by
\begin{equation}
	\left(
	\nu_\pp,
	\nu'_\pp,
	\Omega_\pp
	\right)
	=
	\left(
	-\omega-\nu',
	-\nu',
	\omega+\nu+\nu'
	\right)
	\label{eq:pp_frequency_map}
\end{equation}

These three combinations determine the characteristic frequency structure of
the parquet vertex.
At the lowest non-vanishing order, each reducible vertex depends only on the
bosonic transfer natural to its channel
\begin{align}
	\Phi^{\eh,(2)}
	&\equiv
	\Phi^{\eh,(2)}(\omega)
	\\
	\Phi^{\teh,(2)}
	&\equiv
	\Phi^{\teh,(2)}(\nu'-\nu)
	\\
	\Phi^{\pp,(2)}
	&\equiv
	\Phi^{\pp,(2)}(\nu+\nu'+\omega)
	\label{eq:Phi_frequency_second_order}
\end{align}
Beyond second order, the primitive reducible vertices are genuine
three-frequency objects.
The ordinary channel-pair contribution assigns the natural bosonic transfer
to the central resolvent and the remaining fermionic frequencies to the
coupling vertices. The complete third-order construction also retains the
activated-sector terms identified below.

The exact channel-irreducible vertices satisfy
\begin{align}
	\Gamma^\eh
	&=
	\Lambda
	+
	\Phi^{\teh}
	+
	\Phi^\pp
	\label{eq:Gamma_eh}
	\\
	\Gamma^\pp
	&=
	\Lambda
	+
	\Phi^\eh
	+
	\Phi^{\teh}
	\label{eq:Gamma_pp}
\end{align}
For a two-body interaction, the first correction to the fully irreducible
vertex beyond the bare interaction occurs at fourth order, so that
\begin{equation}
	\Lambda
	=
	\Lambda^{(1)}
	+
	O(v^4)
	\label{eq:Lambda_order}
\end{equation}
where $v$ denotes the electron-electron interaction.
The second-order channel-irreducible vertices are therefore
\begin{align}
	\Gamma^{\eh,(2)}
	&=
	\Phi^{\teh,(2)}
	+
	\Phi^{\pp,(2)}
	\label{eq:Gamma_eh_second}
	\\
	\Gamma^{\pp,(2)}
	&=
	\Phi^{\eh,(2)}
	+
	\Phi^{\teh,(2)}
	\label{eq:Gamma_pp_second}
\end{align}

%%%%%%%%%%%%%%%%%%%%%%%%%
\section{Target multiresolvent representation}
\label{sec:target_representation}
%%%%%%%%%%%%%%%%%%%%%%%%%

Having fixed the channel frequencies, we can state the algebraic form that we
seek before deriving any perturbative ingredient.

For the ordinary channel-pair contribution in each primitive channel, the natural bosonic transfer is carried
by a central resolvent, while the remaining fermionic-frequency dependences
are carried by a left and a right coupling vertex.
We denote these coupling vertices by $U^r$ and $V^r$ (with $r = \eh$ or $\pp$), respectively, and keep
them as independent objects unless a symmetry relation is explicitly used.
We denote this contribution by $\Phi^r_{\mathrm{pair}}$. The target forms are
\begin{align}
	\ii\Phi^\eh_{\mathrm{pair}}(\nu,\nu',\omega)
	&=
	U^\eh(\nu,\omega)
	R^\eh(\omega)
	V^\eh(\nu',\omega)
	\label{eq:target_eh}
	\\
	\ii\Phi^\pp_{\mathrm{pair}}(\nu_\pp,\nu'_\pp,\Omega_\pp)
	&=
	U^\pp(\nu_\pp,\Omega_\pp)
	R^\pp(\Omega_\pp)
	V^\pp(\nu'_\pp,\Omega_\pp)
	\label{eq:target_pp}
\end{align}
with
\begin{align}
	R^\eh(\omega)
	&=
	\left[
	\omega\eta^\eh-M^\eh + \ii 0^+ I
	\right]^{-1}
	\label{eq:target_R_eh}
	\\
	R^\pp(\Omega_\pp)
	&=
	\left[
	\Omega_\pp\eta^\pp-M^\pp + \ii 0^+ I
	\right]^{-1}
	\label{eq:target_R_pp}
\end{align}
Here $\eta^r$ is the channel metric, with eigenvalues $+1$ and $-1$ for
the two propagation sectors and $I$ is the identity operator.
Antisymmetrized two-electron integrals are written directly as
$\mel{pq}{}{rs}$ throughout.

The corresponding particle-number structures of these contributions are
\begin{align}
	\Phi^\eh
	&:\quad
	N\pm1
	\longrightarrow
	N
	\longrightarrow
	N\pm1
	\label{eq:target_sector_eh}
	\\
	\Phi^\pp
	&:\quad
	N\pm1
	\longrightarrow
	N\pm2
	\longrightarrow
	N\pm1
	\label{eq:target_sector_pp}
\end{align}
Thus, the central resolvent distinguishes $\eh$ and $\pp$ propagation,
whereas the coupling vertices retain the fermionic frequencies associated
with $(N\pm1)$-electron propagation.
At second order, the coupling
vertices entering $\Phi^{r,(2)}$ are the first-order quantities
$U^{r,(1)}$ and $V^{r,(1)}$ and are therefore instantaneous.
The non-trivial $\hhp$ and $\pph$ pole structure first appears in
$U^{r,(2)}$ and $V^{r,(2)}$, which enter the third-order reducible vertex and
are evaluated in Appendices~\ref{sec:eh_third} and \ref{sec:pp_third}.

The zeroth-order central resolvent is
\begin{equation}
	R^{r,(0)}(\Omega_r)
	=
	\left[
	\Omega_r \eta^r - M^{r,(0)} + \ii 0^+ I
	\right]^{-1}
	\label{eq:target_R_zero}
\end{equation}
and the second-order target is
\begin{equation}
	\ii\Phi^{r,(2)}
	=
	U^{r,(1)}
	R^{r,(0)}
	V^{r,(1)}
	\label{eq:target_second_order}
\end{equation}

\subsection{Perturbative expansion through third order}
Expanding the ordinary channel-pair contribution by one perturbative order,
\begin{align}
    U^r
    &=
    U^{r,(1)}+U^{r,(2)}+\mathcal O(v^3)
    \\
    V^r
    &=
    V^{r,(1)}+V^{r,(2)}+\mathcal O(v^3)
    \label{eq:target_UV_expansion}
    \\
    M^r
    &=
    M^{r,(0)}+M^{r,(1)}+\mathcal O(v^2)
    \label{eq:target_M_expansion}
\end{align}
then
\begin{align}
    R^r
    &=
    R^{r,(0)}+R^{r,(1)}+\mathcal O(v^2)
    \\
    R^{r,(1)}
    &=
    R^{r,(0)}M^{r,(1)}R^{r,(0)}
    \label{eq:target_R_first}
\end{align}
and one obtains the familiar third-order combination
\begin{equation}
    U^{r,(2)}R^{r,(0)}V^{r,(1)}
    +
    U^{r,(1)}R^{r,(0)}V^{r,(2)}
    +
    U^{r,(1)}R^{r,(1)}V^{r,(1)}
    \label{eq:target_third_order}
\end{equation}
The three terms correspond, respectively, to second-order corrections to the left and right coupling vertices and to a first-order correction to the central resolvent.


This three-term combination describes the ordinary channel-pair contribution.
At second order, $\Phi^r_{\mathrm{pair}}=\Phi^r$; at third order, the complete
vertex also contains activated-sector terms. Their relation to the target
form is established in Sec.~\ref{sec:second_order_multiresolvent}.

\section{Three-frequency BSE and one-sided frequency reduction}
\label{sec:three_frequency_bse}
%%%%%%%%%%%%%%%%%%%%%%%%%

Within the one-shot HF-reference formulation, the $\eh$ correlation function obeys the
four-point BSE
\begin{equation}
	L
	=
	L_0
	+
	L_0 \Gamma^\eh L
	\label{eq:BSE_symbolic}
\end{equation}
In the direct $\eh$ convention of Sec.~\ref{sec:frequencies},
\begin{multline}
	L(\nu,\nu',\omega)
	=
	L_0(\nu,\nu',\omega)
	\\
	+
	\frac{1}{(2\pi)^2}
	\int \dd{\bar\nu}\dd{\bar\nu'}
	L_0(\nu,\bar\nu,\omega)
	\Gamma^\eh(\bar\nu,\bar\nu',\omega)
	L(\bar\nu',\nu',\omega)
	\label{eq:BSE_three_frequency}
\end{multline}
Equation~\eqref{eq:BSE_three_frequency} is simply
Eq.~\eqref{eq:BSE_symbolic} written in the direct-$\eh$ frequency convention
of Sec.~\ref{sec:frequencies}.

To rewrite the BSE as an equation for the full vertex $F$, we use the corresponding relation
\begin{equation}
	L
	=
	L_0
	+
	L_0 F L_0
	\label{eq:L_F_relation}
\end{equation}
Combining Eqs.~\eqref{eq:BSE_symbolic} and \eqref{eq:L_F_relation}, 
and removing the external $L_0$ factors by multiplication with $L_0^{-1}$, gives the vertex BSE
\begin{equation}
	F
	=
	\Gamma^\eh
	+
	\Gamma^\eh L_0^\eh F
	\label{eq:F_BSE_eh}
\end{equation}
Using $F=\Gamma^\eh+\Phi^\eh$ in the direct $\eh$ channel, the reducible
part therefore satisfies
\begin{equation}
	\Phi^\eh
	=
	\Gamma^\eh L_0^\eh F
	\label{eq:Phi_BSE_eh}
\end{equation}
The analogous relation $\Phi^r=\Gamma^rL_0^rF$ holds in each parquet
channel $r$ with the corresponding channel propagation and its symmetry
weight. In the $\pp$ channel, the two orientations of each unique pair
are averaged before the final pair-space contraction.
These reference-propagator vertex relations are expanded perturbatively below.
Barred frequencies denote integration variables throughout these notes.

For occupied and virtual reference orbitals, we have
\begin{align}
	G_{0,i}(\nu)
	&=
	\frac{1}{\nu-\eps_i-\ii\eta}
	\\
	G_{0,a}(\nu)
	&=
	\frac{1}{\nu-\eps_a+\ii\eta}
	\label{eq:G0}
\end{align}
We use $+$ for forward propagation $i\to a$ and $-$ for backward propagation
$a\to i$, while keeping the orbital pair label $ia$ in both cases.
The corresponding three-frequency non-interacting pair propagators are
\begin{align}
	L_{0,ia}^{+}(\nu,\nu',\omega)
	&=
	2\pi\delta(\nu-\nu')
	G_{0,a}(\nu+\omega) G_{0,i}(\nu)
	\label{eq:L0_three_frequency_plus}
	\\
	L_{0,ia}^{-}(\nu,\nu',\omega)
	&=
	2\pi\delta(\nu-\nu')
	G_{0,i}(\nu+\omega)G_{0,a}(\nu)
	\label{eq:L0_three_frequency_minus}
\end{align}

%%%%%%%%%%%%%%%%%%%%%%%%%
\subsection{Exact factorization of the non-interacting pair propagator}
%%%%%%%%%%%%%%%%%%%%%%%%%

The key observation is that each three-frequency pair propagator factorizes
exactly into a two-frequency factor and the usual one-frequency pair
propagator.
Define
\begin{equation}
	\begin{aligned}
	L_{0,ia}^+(\omega)
	&=
	\frac{\ii}
	{\omega-(\eps_a-\eps_i)+\ii0^+}
	\\
	L_{0,ia}^-(\omega)
	&=
	-\frac{\ii}
	{\omega+(\eps_a-\eps_i)-\ii0^+}
	\end{aligned}
	\label{eq:L0_pm}
\end{equation}
and the corresponding one-sided factors
\begin{equation}
	\begin{aligned}
	Q^+_{ia}(\nu,\omega)
	&=
	-\ii
	\qty[
	G_{0,i}(\nu) - G_{0,a}(\nu+\omega)
	]
	\\
	Q^-_{ia}(\nu,\omega)
	&=
	+\ii
	\qty[
	G_{0,a}(\nu) - G_{0,i}(\nu+\omega)
	]
	\end{aligned}
	\label{eq:Q_pm}
\end{equation}
Starting from Eqs.~\eqref{eq:L0_three_frequency_plus} and
\eqref{eq:L0_three_frequency_minus}, substituting the Green's functions of
Eq.~\eqref{eq:G0}, and using the definitions
\eqref{eq:L0_pm} and \eqref{eq:Q_pm}, partial fractioning gives the exact
identities
\begin{align}
	L_{0,ia}^{+}(\nu,\nu',\omega)
	&=
	2\pi\delta(\nu-\nu')
	Q^+_{ia}(\nu,\omega)
	L_{0,ia}^{+}(\omega)
	\label{eq:L0_factorization_plus}
	\\
	L_{0,ia}^{-}(\nu,\nu',\omega)
	&=
	2\pi\delta(\nu-\nu')
	Q^-_{ia}(\nu,\omega)
	L_{0,ia}^{-}(\omega)
	\label{eq:L0_factorization_minus}
\end{align}

%%%%%%%%%%%%%%%%%%%%%%%%%

The one-sided factors are normalized according to
\begin{equation}
    \frac{1}{2\pi}\int\dd{\nu}\,Q^\sigma_{ia}(\nu,\omega)=1,
    \qquad \sigma=\pm.
    \label{eq:Q_normalization}
\end{equation}
The factorization above concerns occupied-virtual and virtual-occupied
propagation. The isolated frequency integrals of the complementary
occupied-occupied and virtual-virtual propagators vanish, but their
convolutions with a dynamic kernel need not vanish.
These additional sectors must therefore be retained at third order, as
shown in Appendix~\ref{sec:eh_third}.

\subsection{One-sided frequency integration}
%%%%%%%%%%%%%%%%%%%%%%%%%

We now insert the exact factorization
\eqref{eq:L0_factorization_plus}--\eqref{eq:L0_factorization_minus} into the
internal non-interacting pair propagation of the vertex BSE
\eqref{eq:Phi_BSE_eh}.
The delta function in the factorized $L_0$ removes one integration variable,
leaving one factor $Q^\sigma$ and one one-frequency propagator
$L_0^\sigma(\omega)$.

Let $\sigma,\sigma'=\pm$ denote the propagation directions of the left pair
$ia$ and right pair $jb$, respectively. Thus $++$ denotes forward propagation
on both pairs, whereas $+-$ denotes forward propagation on $ia$ and backward
propagation on $jb$.
In this pair-space discussion we suppress the channel superscript $\eh$ on
$U$ and $V$.
To match the convention $\ii\Phi^r=U^rR^rV^r$, the projected coupling vertices are defined from $\ii\Gamma^\eh$.

Projecting the left pair gives
\begin{multline}
	\frac{1}{(2\pi)^2}
	\int \dd{\bar\nu_1}\dd{\bar\nu_2}
	L_{0,ia}^{\sigma}(\bar\nu_1,\bar\nu_2,\omega)
	\left[
	\ii\Gamma^\eh
	\right]^{\sigma\sigma'}_{ia,jb}(\bar\nu_2,\nu',\omega)
	\\
	=
	L_{0,ia}^\sigma(\omega)
	\frac{1}{2\pi}
	\int \dd{\bar\nu}
	Q^\sigma_{ia}(\bar\nu,\omega)
	\left[
	\ii\Gamma^\eh
	\right]^{\sigma\sigma'}_{ia,jb}(\bar\nu,\nu',\omega)
	\label{eq:BSE_right_factorization}
\end{multline}
The equality in Eq.~\eqref{eq:BSE_right_factorization} follows directly from
the appropriate factorization identity
\eqref{eq:L0_factorization_plus} or
\eqref{eq:L0_factorization_minus} and the associated delta function.
This motivates the definition of the right coupling vertex
\begin{equation}
	V^{\sigma\sigma'}_{ia,jb}(\nu',\omega)
	=
	\frac{1}{2\pi}
	\int \dd{\bar\nu}
	Q^\sigma_{ia}(\bar\nu,\omega)
	\left[
	\ii\Gamma^\eh
	\right]^{\sigma\sigma'}_{ia,jb}(\bar\nu,\nu',\omega)
	\label{eq:V_def}
\end{equation}
so that Eq.~\eqref{eq:BSE_right_factorization} becomes
\begin{equation}
	L_{0,ia}^\sigma(\omega)
	V^{\sigma\sigma'}_{ia,jb}(\nu',\omega)
	\label{eq:BSE_right_compact}
\end{equation}

Projecting the right pair gives analogously
\begin{multline}
	\frac{1}{(2\pi)^2}
	\int \dd{\bar\nu'_1}\dd{\bar\nu'_2}
	\left[
	\ii\Gamma^\eh
	\right]^{\sigma\sigma'}_{ia,jb}(\nu,\bar\nu'_1,\omega)
	L_{0,jb}^{\sigma'}(\bar\nu'_1,\bar\nu'_2,\omega)
	\\
	=
	\frac{1}{2\pi}
	\int \dd{\bar\nu'}
	\left[
	\ii\Gamma^\eh
	\right]^{\sigma\sigma'}_{ia,jb}(\nu,\bar\nu',\omega)
	Q^{\sigma'}_{jb}(\bar\nu',\omega)
	L_{0,jb}^{\sigma'}(\omega)
	\label{eq:BSE_left_factorization}
\end{multline}
Equation~\eqref{eq:BSE_left_factorization} follows from the same
factorization identities, now applied to the propagator on the right of the
kernel.
This motivates the definition of the left coupling vertex
\begin{equation}
	U^{\sigma\sigma'}_{ia,jb}(\nu,\omega)
	=
	\frac{1}{2\pi}
	\int \dd{\bar\nu'}
	\left[
	\ii\Gamma^\eh
	\right]^{\sigma\sigma'}_{ia,jb}(\nu,\bar\nu',\omega)
	Q^{\sigma'}_{jb}(\bar\nu',\omega)
	\label{eq:U_def}
\end{equation}
and therefore
\begin{equation}
	U^{\sigma\sigma'}_{ia,jb}(\nu,\omega)
	L_{0,jb}^{\sigma'}(\omega)
	\label{eq:BSE_left_compact}
\end{equation}

Thus $U$ retains the left fermionic frequency $\nu$, whereas $V$ retains the
right fermionic frequency $\nu'$.
Equations~\eqref{eq:U_def} and \eqref{eq:V_def} are therefore definitions
of exact one-sided projections of the three-frequency irreducible vertex.
They use only the exact factorization of $L_0$,
Eqs.~\eqref{eq:L0_factorization_plus} and
\eqref{eq:L0_factorization_minus}.
Perturbative approximations enter only when $\Gamma$, the coupling vertices,
and the central resolvent are truncated in powers of the interaction.

%%%%%%%%%%%%%%%%%%%%%%%%%
\subsection{From three to two to one frequency}
%%%%%%%%%%%%%%%%%%%%%%%%%

This second projection is included because it is used
in the independent GF2 consistency check summarized later.
Starting from the exact one-sided definitions
\eqref{eq:U_def} and \eqref{eq:V_def}, a second integration removes the
remaining fermionic frequency.
We define the resulting one-frequency projected $\eh$ kernel by
\begin{multline}
	\left[
	\ii\widetilde{\Gamma}^\eh
	\right]^{\sigma\sigma'}_{ia,jb}(\omega)
	=
	\frac{1}{(2\pi)^2}
	\int \dd{\bar\nu}\dd{\bar\nu'}
	Q^\sigma_{ia}(\bar\nu,\omega)
	\\
	\times
	\left[
	\ii\Gamma^\eh
	\right]^{\sigma\sigma'}_{ia,jb}(\bar\nu,\bar\nu',\omega)
	Q^{\sigma'}_{jb}(\bar\nu',\omega)
	\label{eq:Gamma_tilde_def}
\end{multline}
Substituting Eq.~\eqref{eq:V_def} or, equivalently,
Eq.~\eqref{eq:U_def} into Eq.~\eqref{eq:Gamma_tilde_def}, and merely
reordering the two integrations, gives the same object in either order
\begin{align}
	\left[
	\ii\widetilde{\Gamma}^\eh
	\right]^{\sigma\sigma'}_{ia,jb}(\omega)
	&=
	\frac{1}{2\pi}
	\int \dd{\bar\nu'}
	V^{\sigma\sigma'}_{ia,jb}(\bar\nu',\omega)
	Q^{\sigma'}_{jb}(\bar\nu',\omega)
	\label{eq:Gamma_tilde_from_V}
	\\
	&=
	\frac{1}{2\pi}
	\int \dd{\bar\nu}
	Q^\sigma_{ia}(\bar\nu,\omega)
	U^{\sigma\sigma'}_{ia,jb}(\bar\nu,\omega)
	\label{eq:Gamma_tilde_from_U}
\end{align}
The frequency reduction is therefore
\begin{equation}
	\Gamma^\eh(\nu,\nu',\omega)
	\longrightarrow
	\left\{
	U(\nu,\omega),V(\nu',\omega)
	\right\}
	\longrightarrow
	\widetilde{\Gamma}^\eh(\omega)
	\label{eq:frequency_reduction}
\end{equation}
The first arrow denotes one-sided frequency integration and the second the
integration over the remaining fermionic frequency.

%%%%%%%%%%%%%%%%%%%%%%%%%
\section{Second-order primitive vertices and crossing symmetry}
\label{sec:second_order_primitives}
%%%%%%%%%%%%%%%%%%%%%%%%%

We now construct the primitive reducible vertices at their lowest
non-vanishing order and then combine them with the parquet relations of
Sec.~\ref{sec:frequencies}.
The only perturbative ingredients needed are the bare irreducible interaction
$\Gamma^{(1)}$ and the non-interacting two-particle propagation $L_0$.

%%%%%%%%%%%%%%%%%%%%%%%%%
\subsection{Direct eh reducible vertex}
%%%%%%%%%%%%%%%%%%%%%%%%%

At first order, the $\eh$ irreducible interaction is the bare
antisymmetrized interaction
\begin{equation}
	\ii\Gamma^{\eh,(1)}_{pqrs}
	=
	\mel{pq}{}{rs}
	\label{eq:Gamma_eh_first}
\end{equation}
We now expand the exact vertex BSE Eq.~\eqref{eq:Phi_BSE_eh} in powers of
the bare interaction.
At first order, the full and irreducible vertices coincide,
$F^{(1)}=\Gamma^{\eh,(1)}$, because the only contribution is the bare
interaction of Eq.~\eqref{eq:Gamma_eh_first}.
Keeping the second-order term of Eq.~\eqref{eq:Phi_BSE_eh} therefore gives
\begin{equation}
	\begin{aligned}
	\ii\Phi^{\eh,(2)}
	&=
	\left[\ii\Gamma^{\eh,(1)}\right]
	\circ\left[-\ii L_0^\eh\right]
	\circ\left[\ii F^{(1)}\right]
	\\
	&=
	\left[\ii\Gamma^{\eh,(1)}\right]
	\circ\left[-\ii L_0^\eh\right]
	\circ\left[\ii\Gamma^{\eh,(1)}\right]
	\end{aligned}
	\label{eq:Phi_eh_second_convolution}
\end{equation}
Equation~\eqref{eq:Phi_eh_second_convolution}
is simply the second-order component of the exact vertex BSE.
Using Eq.~\eqref{eq:Gamma_eh_first} for the two end vertices and
Eqs.~\eqref{eq:L0_three_frequency_plus},
\eqref{eq:L0_three_frequency_minus}, and \eqref{eq:L0_pm} for the
non-interacting $\eh$ propagation gives
\begin{equation}
	\begin{split}
	\ii\Phi^{\eh,(2)}_{pqrs}(\omega)
	&=
	\sum_{ia}
	\frac{
	\mel{pa}{}{ri}
	\mel{qi}{}{sa}
	}
	{\omega-(\eps_a-\eps_i)+\ii0^+}
	\\
	&-
	\sum_{ia}
	\frac{
	\mel{pi}{}{ra}
	\mel{qa}{}{si}
	}
	{\omega+(\eps_a-\eps_i)-\ii0^+}
	\end{split}
	\label{eq:Phi_eh_second}
\end{equation}
At this order the two end vertices are instantaneous and
$\Phi^{\eh,(2)}$ depends only on the direct $\eh$ transfer
frequency $\omega$.

%%%%%%%%%%%%%%%%%%%%%%%%%
\subsection{Transverse eh contribution from crossing}\label{sec:crossing}
%%%%%%%%%%%%%%%%%%%%%%%%%

Fermionic crossing generates the transverse contribution order by order
\begin{equation}
	\Phi^{\teh,(n)}_{pqrs}(\nu,\nu',\omega)
	=
	-
	\Phi^{\eh,(n)}_{pqsr}
	\left(
	\nu,
	\nu+\omega,
	\nu'-\nu
	\right)
	\label{eq:Phi_ehbar_crossing}
\end{equation}
Applying the crossing relation
Eq.~\eqref{eq:Phi_ehbar_crossing} to the direct contribution
Eq.~\eqref{eq:Phi_eh_second}, with the crossed transfer
$\Omega_{\teh}=\nu'-\nu$ from Eq.~\eqref{eq:Omega_ehbar}, yields
\begin{equation}
	\begin{split}
	\ii\Phi^{\teh,(2)}_{pqrs}(\nu,\nu',\omega)
	&=
	-
	\sum_{kc}
	\frac{
	\mel{pc}{}{sk}
	\mel{qk}{}{rc}
	}
	{\nu'-\nu-(\eps_c-\eps_k)+\ii0^+}
	\\
	&+
	\sum_{kc}
	\frac{
	\mel{pk}{}{sc}
	\mel{qc}{}{rk}
	}
	{\nu'-\nu+(\eps_c-\eps_k)-\ii0^+}
	\end{split}
	\label{eq:Phi_ehbar_second}
\end{equation}
The characteristic transverse $\eh$ poles are therefore controlled
by $\nu'-\nu$.

At third order, the same crossing relation is applied to the complete
direct-$\eh$ result, including the ordinary channel-pair corrections,
the activated $\text{oo}/\text{vv}$ contributions, and the three-interaction
term. The direct result is summarized in
Sec.~\ref{sec:second_order_multiresolvent} and derived in
Appendix~\ref{sec:eh_third}.
The transverse channel requires no independent third-order convolution;
all its terms follow by crossing the complete direct contribution.

%%%%%%%%%%%%%%%%%%%%%%%%%
\subsection{pp reducible vertex}
%%%%%%%%%%%%%%%%%%%%%%%%%

The $\pp$ channel contains $\pp$ and $\hh$ propagation.
Applying the channel analogue of Eq.~\eqref{eq:Phi_BSE_eh},
$\Phi^\pp=\Gamma^\pp L_0^\pp F$, and retaining its second-order component
produces the same two-bare-vertex structure as
Eq.~\eqref{eq:Phi_eh_second_convolution}, now in the native transfer
frequency $\Omega_\pp$ defined by Eq.~\eqref{eq:Omega_pp}.
This gives
\begin{equation}
	\begin{split}
	\ii\Phi_{pqrs}^{\pp,(2)}(\Omega)
	&=
	\sum_{c<d}
	\frac{
	\mel{pq}{}{cd}
	\mel{cd}{}{rs}
	}
	{\Omega-(\eps_c+\eps_d)+\ii0^+}
	\\
	&-
	\sum_{k<l}
	\frac{
	\mel{pq}{}{kl}
	\mel{kl}{}{rs}
	}
	{\Omega-(\eps_k+\eps_l)-\ii0^+}
	\end{split}
	\label{eq:Phi_pp_second_native}
\end{equation}

To combine this contribution with the direct-$\eh$ parquet relation, we
rewrite it in the direct-$\eh$ variables.
Substituting the frequency map Eq.~\eqref{eq:pp_frequency_map}, in
particular $\Omega_\pp=\nu+\nu'+\omega$, into
Eq.~\eqref{eq:Phi_pp_second_native} gives
\begin{equation}
	\begin{split}
	\ii\Phi^{\pp,(2)}_{pqrs}(\nu,\nu',\omega)
	&=
	\sum_{c<d}
	\frac{
	\mel{pq}{}{cd}
	\mel{cd}{}{rs}
	}
	{\nu+\nu'+\omega-\eps_c-\eps_d+\ii0^+}
	\\
	&-
	\sum_{k<l}
	\frac{
	\mel{pq}{}{kl}
	\mel{kl}{}{rs}
	}
	{\nu+\nu'+\omega-\eps_k-\eps_l-\ii0^+}
	\end{split}
	\label{eq:Phi_pp_second_eh}
\end{equation}
The characteristic pair poles are therefore controlled by
$\nu+\nu'+\omega$.

%%%%%%%%%%%%%%%%%%%%%%%%%
\subsection{Second-order eh irreducible kernel}
%%%%%%%%%%%%%%%%%%%%%%%%%

At second order, Eq.~\eqref{eq:Lambda_order} implies that the fully
irreducible vertex contributes only through the bare first-order
interaction.
Therefore the direct-$\eh$ parquet relation
Eq.~\eqref{eq:Gamma_eh_second} contains only the transverse-$\eh$ and
$\pp$ reducible pieces,
\begin{equation}
	\ii\Gamma^{\eh,(2)}
	=
	\ii\Phi^{\teh,(2)}
	+
	\ii\Phi^{\pp,(2)}
	\label{eq:Gamma_eh_second_i}
\end{equation}
Substituting Eqs.~\eqref{eq:Phi_ehbar_second} and
\eqref{eq:Phi_pp_second_eh} into Eq.~\eqref{eq:Gamma_eh_second_i} gives
explicitly
\begin{equation}
	\begin{split}
	\ii\Gamma^{\eh,(2)}_{pqrs}(\nu,\nu',\omega)
	&=
	-
	\sum_{kc}
	\frac{
	\mel{pc}{}{sk}
	\mel{qk}{}{rc}
	}
	{\nu'-\nu-(\eps_c-\eps_k)+\ii0^+}
	\\
	&+
	\sum_{kc}
	\frac{
	\mel{pk}{}{sc}
	\mel{qc}{}{rk}
	}
	{\nu'-\nu+(\eps_c-\eps_k)-\ii0^+}
	\\
	&+
	\sum_{c<d}
	\frac{
	\mel{pq}{}{cd}
	\mel{cd}{}{rs}
	}
	{\nu+\nu'+\omega-\eps_c-\eps_d+\ii0^+}
	\\
	&-
	\sum_{k<l}
	\frac{
	\mel{pq}{}{kl}
	\mel{kl}{}{rs}
	}
	{\nu+\nu'+\omega-\eps_k-\eps_l-\ii0^+}
	\end{split}
	\label{eq:Gamma_eh_second_explicit}
\end{equation}
Equation~\eqref{eq:Gamma_eh_second_explicit} is the complete
second-order direct-$\eh$ irreducible kernel.
Its one-sided projections generate the second-order coupling vertices needed
at third order; those explicit projections are given in Appendix~\ref{sec:eh_third}.

%%%%%%%%%%%%%%%%%%%%%%%%%
\section{Multiresolvent construction through third order}
\label{sec:second_order_multiresolvent}
%%%%%%%%%%%%%%%%%%%%%%%%%

\subsection{Second-order construction}

We now verify the target representation
Eq.~\eqref{eq:target_second_order} directly against the explicit primitive
vertices derived above.
At this order, the end vertices are the bare interaction
Eq.~\eqref{eq:Gamma_eh_first}, while the central propagation is the
zeroth-order non-interacting resolvent.

For the $\eh$ channel, the two denominators in
Eq.~\eqref{eq:Phi_eh_second} are reproduced by using the general resolvent
definition Eq.~\eqref{eq:target_R_zero} with
\begin{align}
	\eta^\eh
	&=
	\begin{pmatrix}
	1&0
	\\
	0&-1
	\end{pmatrix}
	&
	M^{\eh,(0)}
	&=
	\begin{pmatrix}
	\Delta\eps&0
	\\
	0&\Delta\eps
	\end{pmatrix}
	\label{eq:eta_M_eh_zero}
\end{align}
with
\begin{equation}
	(\Delta\eps)_{ia,jb}
	=
	(\eps_a-\eps_i)\delta_{ij}\delta_{ab}
	\label{eq:Delta_eh}
\end{equation}
Using Eqs.~\eqref{eq:eta_M_eh_zero} and \eqref{eq:Delta_eh} in
Eq.~\eqref{eq:target_R_zero} gives
\begin{equation}
	R^{\eh,(0)}(\omega)
	=
	\left[
	\omega\eta^\eh-M^{\eh,(0)}+\ii0^+I
	\right]^{-1}
	\label{eq:R_eh_zero}
\end{equation}
The two blocks of Eq.~\eqref{eq:R_eh_zero} are exactly the forward ($+$)
and backward ($-$) denominators appearing in
Eq.~\eqref{eq:Phi_eh_second}.
Since the end vertices in Eq.~\eqref{eq:Phi_eh_second_convolution} are the
bare interaction Eq.~\eqref{eq:Gamma_eh_first}, the corresponding
first-order coupling-vertex components are
\begin{align}
	U^{\eh,+,(1)}_{pr,ia}
	&=
	\mel{pa}{}{ri}
	&
	V^{\eh,+,(1)}_{ia,sq}
	&=
	\mel{qi}{}{sa}
	\label{eq:UV_eh_first_plus}
	\\
	U^{\eh,-,(1)}_{pr,ia}
	&=
	\mel{pi}{}{ra}
	&
	V^{\eh,-,(1)}_{ia,sq}
	&=
	\mel{qa}{}{si}
	\label{eq:UV_eh_first_minus}
\end{align}
Substituting Eqs.~\eqref{eq:R_eh_zero},
\eqref{eq:UV_eh_first_plus}, and \eqref{eq:UV_eh_first_minus} into the
target Eq.~\eqref{eq:target_second_order} gives
\begin{equation}
	\ii\Phi^{\eh,(2)}
	=
	U^{\eh,(1)}R^{\eh,(0)}V^{\eh,(1)}
	\label{eq:Phi_eh_second_resolvent}
\end{equation}
which reproduces Eq.~\eqref{eq:Phi_eh_second} term by term.

For the $\pp$ channel, we proceed in exactly the same way.
The particle-pair and hole-pair denominators in
Eq.~\eqref{eq:Phi_pp_second_native} are generated from
Eq.~\eqref{eq:target_R_zero} by choosing
\begin{align}
	\eta^\pp
	&=
	\begin{pmatrix}
	1&0
	\\
	0&-1
	\end{pmatrix}
	&
	M^{\pp,(0)}
	&=
	\begin{pmatrix}
	\eps_a+\eps_b&0
	\\
	0&-(\eps_i+\eps_j)
	\end{pmatrix}
	\label{eq:eta_M_pp_zero}
\end{align}
with $a<b$ and $i<j$.
Inserting Eq.~\eqref{eq:eta_M_pp_zero} into
Eq.~\eqref{eq:target_R_zero} gives
\begin{equation}
	R^{\pp,(0)}(\Omega_\pp)
	=
	\left[
	\Omega_\pp\eta^\pp-M^{\pp,(0)}+\ii0^+I
	\right]^{-1}
	\label{eq:R_pp_zero}
\end{equation}
The end vertices are again the bare interaction, so the first-order
coupling vertices are
\begin{align}
	U^{\pp,(1)}_{pq,ab}
	&=
	\mel{pq}{}{ab}
	&
	V^{\pp,(1)}_{ab,rs}
	&=
	\mel{ab}{}{rs}
	\label{eq:UV_pp_first_ee}
	\\
	U^{\pp,(1)}_{pq,ij}
	&=
	\mel{pq}{}{ij}
	&
	V^{\pp,(1)}_{ij,rs}
	&=
	\mel{ij}{}{rs}
	\label{eq:UV_pp_first_hh}
\end{align}
Substituting Eqs.~\eqref{eq:R_pp_zero},
\eqref{eq:UV_pp_first_ee}, and \eqref{eq:UV_pp_first_hh} into the target
Eq.~\eqref{eq:target_second_order} therefore gives
\begin{equation}
	\ii\Phi^{\pp,(2)}
	=
	U^{\pp,(1)}R^{\pp,(0)}V^{\pp,(1)}
	\label{eq:Phi_pp_second_resolvent}
\end{equation}
which reproduces Eq.~\eqref{eq:Phi_pp_second_native}.
The channel metric supplies the relative sign of the two sectors,
while the $+\ii0^+I$ prescription in the inverse resolvent
produces their opposite causal pole shifts.

Equations~\eqref{eq:Phi_eh_second_resolvent} and
\eqref{eq:Phi_pp_second_resolvent} establish the primitive multiresolvent
representation completely through second order.



%%%%%%%%%%%%%%%%%%%%%%%%%

\subsection{Complete third-order construction}

At third order, the ordinary channel-pair contribution contains second-order
corrections to the left and right couplings and a first-order correction
to central propagation, as in Eq.~\eqref{eq:target_third_order}.
The first-order interaction matrices $M^{\eh,(1)}$ and $M^{\pp,(1)}$
are given in Eqs.~\eqref{eq:M_eh_first} and \eqref{eq:M_pp_first}.
The explicit $U^{r,(2)}$ and $V^{r,(2)}$ expressions in
Appendices~\ref{sec:eh_third} and \ref{sec:pp_third} contain
$\hhp$ and $\pph$ denominators and retain one external fermionic frequency.

The frequency-dependent second-order irreducible kernel also activates
orbital sectors whose isolated reference-propagator frequency integrals
vanish. Their role is complementary in the two primitive channels:
\begin{table}[b]
\caption{Internal orbital sectors in the third-order primitive vertices.
The activated sectors contribute when the second-order irreducible kernel
is frequency dependent.}
\label{tab:orbital_sectors}
\begin{ruledtabular}
\begin{tabular}{ccc}
Channel & Ordinary pair sectors & Activated sectors \\
$\eh$ & $\text{ov},\text{vo}$ & $\text{oo},\text{vv}$ \\
$\pp$ & $\text{oo},\text{vv}$ & $\text{ov},\text{vo}$
\end{tabular}
\end{ruledtabular}
\end{table}

Let $\mathcal L_0^{r,x}$ denote the reduced reference propagation in orbital
sector $x$, with the phase and pair-orientation weights appropriate to the
channel. These propagators are defined in
Eqs.~\eqref{eq:L0_eh_reduced_sectors} and \eqref{eq:L0_pp_reduced_oriented}.
The additional third-order contributions are
\begin{equation}
\begin{split}
\mathcal E^{r,(3)}
&=\sum_{x\in\mathcal A_r}\Big(
[\ii\Gamma^{r,(2)}]\circ\mathcal L_0^{r,x}
\circ[\ii\Gamma^{r,(1)}]
\\
&\qquad+[\ii\Gamma^{r,(1)}]\circ\mathcal L_0^{r,x}
\circ[\ii\Gamma^{r,(2)}]\Big),
\end{split}
\label{eq:activated_third}
\end{equation}
where $\mathcal A_\eh=\{\text{oo},\text{vv}\}$ and
$\mathcal A_\pp=\{\text{ov},\text{vo}\}$.
The complete result is therefore
\begin{equation}
\begin{split}
\ii\Phi^{r,(3)}
&=U^{r,(2)}R^{r,(0)}V^{r,(1)}
+U^{r,(1)}R^{r,(0)}V^{r,(2)}
\\
&\quad+U^{r,(1)}R^{r,(1)}V^{r,(1)}
+\mathcal E^{r,(3)}.
\end{split}
\label{eq:complete_third_main}
\end{equation}
Here $\mathcal E^{r,(3)}$ is already the contribution to $\ii\Phi^{r,(3)}$;
no further phase factor is implied.
For the $\eh$ channel, its explicit terms are
Eqs.~\eqref{eq:Phi_eh_third_vv_left}--\eqref{eq:Phi_eh_third_oo_right};
for the $\pp$ channel, they are
Eqs.~\eqref{eq:Phi_pp_third_mixed_left} and
\eqref{eq:Phi_pp_third_mixed_right}.
They are expressed as products of charged-state denominators without
introducing artificial differences between same-occupancy orbital energies.
The full orbital-sector expressions are collected in
Eqs.~\eqref{eq:Phi_eh_third_final} and \eqref{eq:Phi_pp_third_final}.

Equations~\eqref{eq:target_second_order} and
\eqref{eq:complete_third_main} specify the primitive vertices through third
order. The additional sector terms are part of the complete result; the
three-term expansion of the ordinary $URV$ product alone describes only
the ordinary channel-pair contribution.

\section{Channel consistency and validation}
\label{sec:channel_consistency}
%%%%%%%%%%%%%%%%%%%%%%%%%

The second- and third-order constructions admit complementary consistency checks.
\subsection{Frequency maps and second-order reductions}
The first check is crossing: applying
Eq.~\eqref{eq:Phi_ehbar_crossing} to
Eq.~\eqref{eq:Phi_eh_second} must reproduce
Eq.~\eqref{eq:Phi_ehbar_second}.
This simultaneously checks the orbital permutation and the transverse
transfer $\nu'-\nu$ of Eq.~\eqref{eq:Omega_ehbar}.

The second check uses the $\pp$ convention: substituting the exact map
Eq.~\eqref{eq:pp_frequency_map} into the native expression
Eq.~\eqref{eq:Phi_pp_second_native} must reproduce
Eq.~\eqref{eq:Phi_pp_second_eh}, with the transfer
$\Omega_\pp=\nu+\nu'+\omega$ fixed by Eq.~\eqref{eq:Omega_pp}.

Additional, deliberately redundant checks are obtained by completing the
second frequency integration.
In the $\eh$ channel, the resulting single-frequency second-order kernel
coincides with the correlation part of the GF2 BSE kernel.
In the $\pp$ channel, using the complete antisymmetrized non-interacting
$\pp$ propagator recovers all three independent GF2 blocks.\cite{MarieRomanielloBlaseLoos2025}
Applying the same frequency-reduction procedure to the conventional dynamical
$GW$ kernel at first order in the screened interaction $W$ likewise recovers
the known dynamical $\eh$ and $\pp$ BSE kernels.\cite{MarieAmmarLoos2024,MarieRomanielloLoos2024,MarieRomanielloBlaseLoos2025}
These calculations do not introduce any new element into the multiresolvent
construction and are therefore not reproduced here.
Their role is purely diagnostic.


\subsection{Orbital-sector completeness}
%%%%%%%%%%%%%%%%%%%%%%%%%

The decompositions in Eqs.~\eqref{eq:L0_sector_decomposition} and
\eqref{eq:L0_pp_sector_decomposition} are identities at the level of the
complete internal orbital sum.
In the $\eh$ channel, the instantaneous kernel selects $\text{ov}/\text{vo}$ propagation
and the dynamic second-order kernel additionally activates $\text{oo}/\text{vv}$.
In the $\pp$ channel, the instantaneous kernel selects $\text{oo}/\text{vv}$ propagation
and the dynamic second-order kernel additionally activates $\text{ov}/\text{vo}$.
Thus all four sectors $\text{oo}$, $\text{ov}$, $\text{vo}$, and $\text{vv}$ have been accounted for in
both primitive channels.

%%%%%%%%%%%%%%%%%%%%%%%%%
\subsection{Electron-hole pair reversal}
%%%%%%%%%%%%%%%%%%%%%%%%%

The backward $\eh$ coupling vertices follow from
\begin{equation}
    Q^-_{ia}(\nu,\omega)
    =
    Q^+_{ia}(\nu+\omega,-\omega)
    \label{eq:Q_reversal_third}
\end{equation}
together with the corresponding orbital-index reversal.
For real orbitals,
\begin{align}
    U^{\eh,-,(2)}_{pr,ia}(\nu,\omega)
    &=
    U^{\eh,+,(2)}_{rp,ia}(\nu+\omega,-\omega)
    \label{eq:U_eh_pair_reversal}
    \\
    V^{\eh,-,(2)}_{ia,sq}(\nu',\omega)
    &=
    V^{\eh,+,(2)}_{ia,qs}(\nu'+\omega,-\omega)
    \label{eq:V_eh_pair_reversal}
\end{align}
Direct substitution reproduces the explicit backward expressions, including
their causal prescriptions.

%%%%%%%%%%%%%%%%%%%%%%%%%
\subsection{Particle-particle antisymmetry}
%%%%%%%%%%%%%%%%%%%%%%%%%

The native second-order $\pp$ kernel obeys the pair-exchange identities
\begin{align}
    \ii\Gamma_{pqsr}^{\pp,(2)}(x,-\Omega-y,\Omega)
    &=
    -\ii\Gamma_{pqrs}^{\pp,(2)}(x,y,\Omega)
    \\
    \ii\Gamma_{qprs}^{\pp,(2)}(-\Omega-x,y,\Omega)
    &=
    -\ii\Gamma_{pqrs}^{\pp,(2)}(x,y,\Omega)
    \label{eq:Gamma_pp_pair_exchange}
\end{align}
which combine orbital exchange with the corresponding reversal of the native
relative frequency.
The explicit coupling vertices satisfy
\begin{align}
    U_{pq,ba}^{\pp,(2)}&=-U_{pq,ab}^{\pp,(2)}
    &
    U_{pq,ji}^{\pp,(2)}&=-U_{pq,ij}^{\pp,(2)}
    \\
    V_{ba,rs}^{\pp,(2)}&=-V_{ab,rs}^{\pp,(2)}
    &
    V_{ji,rs}^{\pp,(2)}&=-V_{ij,rs}^{\pp,(2)}
    \label{eq:UV_pp_antisymmetry}
\end{align}
The factor $1/2$ associated with summing both pair orientations is canceled
when the two mixed orientations $ia$ and $ai$ are combined, giving the
normalization in Eqs.~\eqref{eq:Phi_pp_third_mixed_left} and
\eqref{eq:Phi_pp_third_mixed_right}.

%%%%%%%%%%%%%%%%%%%%%%%%%

\subsection{Crossing and parquet reconstruction through third order}

The transverse contribution follows from the complete direct-$\eh$ result
by Eq.~\eqref{eq:Phi_ehbar_crossing}, as discussed in
Sec.~\ref{sec:crossing}. At third order this includes the crossed images
of all four orbital sectors and the three-interaction term.
Because $\Lambda$ has no second- or third-order correction, the reconstruction
is identical in every channel. For $n=2,3$,
\begin{align}
F^{(n)}&=\Phi^{\eh,(n)}+\Phi^{\teh,(n)}+\Phi^{\pp,(n)},
\\
\Gamma^{\eh,(n)}&=\Phi^{\teh,(n)}+\Phi^{\pp,(n)},
\\
\Gamma^{\pp,(n)}&=\Phi^{\eh,(n)}+\Phi^{\teh,(n)},
\label{eq:parquet_order_reconstruction}
\\
F^{(n)}&=\Gamma^{\eh,(n)}+\Phi^{\eh,(n)}
=\Gamma^{\pp,(n)}+\Phi^{\pp,(n)}.
\label{eq:channel_consistency_through_third}
\end{align}
The same channel organization reconstructs the direct $\eh$ Hedin kernel
through Eq.~\eqref{eq:Gamma_eh}.


\appendix

\section{Direct eh vertex through third order}
\label{sec:eh_third}
%%%%%%%%%%%%%%%%%%%%%%%%%

The direct-$\eh$ construction uses without repetition the second-order
irreducible kernel $\Gamma^{\eh,(2)}$, the normalized factors $Q^\pm$,
the zeroth-order neutral resolvent $R^{\eh,(0)}$, and the first-order
coupling vertices $U^{\eh,(1)}$ and $V^{\eh,(1)}$ defined in Sec.~\ref{sec:second_order_multiresolvent}.
The purpose of this section is to evaluate every orbital sector of the
third-order frequency convolution before reorganizing the $\text{ov}/\text{vo}$ subset in
multiresolvent form.

%%%%%%%%%%%%%%%%%%%%%%%%%
\subsection{Perturbative BSE expansion}
%%%%%%%%%%%%%%%%%%%%%%%%%

At third order, the direct-$\eh$ reducible vertex contains three
contributions.
Suppressing the explicit internal-frequency integrations with the convolution
symbol $\circ$, we absorb the phase associated with the internal pair
propagation into the reduced propagator $\mathcal L_0^\eh$, defined below.
The third-order expansion is then
\begin{equation}
    \begin{split}
    \ii\Phi^{\eh,(3)}
    & =
    \left[\ii\Gamma^{\eh,(2)}\right]
    \circ\mathcal L_0^\eh
    \circ\left[\ii\Gamma^{\eh,(1)}\right]
    \\
    &+
    \left[\ii\Gamma^{\eh,(1)}\right]
    \circ\mathcal L_0^\eh
    \circ\left[\ii\Gamma^{\eh,(2)}\right]
    \\
    &+
    \left[\ii\Gamma^{\eh,(1)}\right]
    \circ\mathcal L_0^\eh
    \circ\left[\ii\Gamma^{\eh,(1)}\right]
    \circ\mathcal L_0^\eh
    \circ\left[\ii\Gamma^{\eh,(1)}\right]
    \end{split}
    \label{eq:Phi_eh_third_convolution}
\end{equation}
The first two terms contain the frequency-dependent second-order irreducible
kernel on opposite sides of the internal propagation.
The last term contains three instantaneous first-order interactions.

For a two-body interaction, the fully two-particle-irreducible vertex has no
second- or third-order contribution beyond the bare interaction,
\begin{equation}
    \ii\Lambda^{(1)}_{pqrs}
    =
    \mel{pq}{}{rs}
    \qquad
    \Lambda
    =
    \Lambda^{(1)}+\mathcal O(v^4)
    \label{eq:Lambda_through_third}
\end{equation}
so that Eq.~\eqref{eq:Phi_eh_third_convolution} contains the complete
third-order direct-$\eh$ reducible contribution.

%%%%%%%%%%%%%%%%%%%%%%%%%
\subsection{Orbital decomposition of the internal propagation}
%%%%%%%%%%%%%%%%%%%%%%%%%

We resolve the reduced internal pair propagator as
\begin{equation}
    \mathcal L_0^\eh
    =
    \mathcal L_0^{\eh,\text{oo}}
    +
    \mathcal L_0^{\eh,\text{ov}}
    +
    \mathcal L_0^{\eh,\text{vo}}
    +
    \mathcal L_0^{\eh,\text{vv}}
    \label{eq:L0_sector_decomposition}
\end{equation}
where $\text{oo}$ and $\text{vv}$ contain the $ij$ and $ab$ sectors, respectively,
while $\text{ov}$ contains the $ia$ sector associated with $Q^+$ and $\text{vo}$ contains
the $ai$ sector associated with $Q^-$.  Explicitly,
\begin{align}
    \mathcal L_{0,ia}^{\eh}(x,\omega)
    &=
    -\ii G_{0,a}(x+\omega)G_{0,i}(x)
    \\
    \mathcal L_{0,ai}^{\eh}(x,\omega)
    &=
    -\ii G_{0,i}(x+\omega)G_{0,a}(x)
    \\
    \mathcal L_{0,ij}^{\eh}(x,\omega)
    &=
    -\ii G_{0,i}(x+\omega)G_{0,j}(x)
    \\
    \mathcal L_{0,ab}^{\eh}(x,\omega)
    &=
    -\ii G_{0,a}(x+\omega)G_{0,b}(x)
    \label{eq:L0_eh_reduced_sectors}
\end{align}
Their isolated frequency integrals are
\begin{align}
    \frac{1}{2\pi}\int\dd{x}\,
    \mathcal L_{0,ia}^{\eh}(x,\omega)
    &=
    \frac{1}{\omega-(\eps_a-\eps_i)+\ii0^+}
    \label{eq:eh_ov_integral}
    \\
    \frac{1}{2\pi}\int\dd{x}\,
    \mathcal L_{0,ai}^{\eh}(x,\omega)
    &=
    -\frac{1}{\omega+(\eps_a-\eps_i)-\ii0^+}
    \label{eq:eh_vo_integral}
    \\
    \frac{1}{2\pi}\int\dd{x}\,
    \mathcal L_{0,ij}^{\eh}(x,\omega)
    &=0
    \\
    \frac{1}{2\pi}\int\dd{x}\,
    \mathcal L_{0,ab}^{\eh}(x,\omega)
    &=0
    \label{eq:same_occupancy_bubbles_zero}
\end{align}
Thus, the $\text{ov}$ and $\text{vo}$ sectors generate the ordinary forward and
backward blocks of the neutral resolvent $R^{\eh,(0)}$, whereas the isolated
$\text{oo}$ and $\text{vv}$ sectors vanish because their two one-particle poles lie in
the same half-plane.  At third order, however, $\Gamma^{\eh,(2)}$ supplies an
additional pole and activates these same-occupancy sectors.

The two elementary contour identities needed for the activated sectors are
\begin{align}
    \frac{1}{2\pi}
    \int\dd{x}
    \frac{\mathcal L_{0,ab}^{\eh}(x,\omega)}{x-A-\ii0^+}
    & =
    \frac{1}
    {(A+\omega-\eps_a+\ii0^+)(A-\eps_b+\ii0^+)}
    \label{eq:vv_elementary_contour}
    \\
    \frac{1}{2\pi}
    \int\dd{x}
    \frac{\mathcal L_{0,ij}^{\eh}(x,\omega)}{x-A+\ii0^+}
    & =
    -\frac{1}
    {(A+\omega-\eps_i-\ii0^+)(A-\eps_j-\ii0^+)}
    \label{eq:oo_elementary_contour}
\end{align}
The same integrals vanish when the causal prescription of the last factor is
reversed.  Thus, the $\text{vv}$ sector selects poles of the second-order kernel
in the upper half-plane, whereas the $\text{oo}$ sector selects poles in the
lower half-plane.  The third term of
Eq.~\eqref{eq:Phi_eh_third_convolution} contains only instantaneous vertices.
Its $\text{oo}$ and $\text{vv}$ integrations therefore vanish, so both internal
propagations in that term belong to the $\text{ov}/\text{vo}$ sectors.

%%%%%%%%%%%%%%%%%%%%%%%%%
\subsection{$\text{ov}$ and $\text{vo}$ sectors}
%%%%%%%%%%%%%%%%%%%%%%%%%

The $\text{ov}$ and $\text{vo}$ parts of the first two terms in
Eq.~\eqref{eq:Phi_eh_third_convolution} can be reorganized as corrections to
the coupling vertices.
For the $\text{ov}$ sector, the left correction is
\begin{equation}
    U^{\eh,+,(2)}_{pr,ia}(\nu,\omega)
    =
    \frac{1}{2\pi}
    \int\dd{\bar\nu}
    \left[
    \ii\Gamma^{\eh,(2)}_{pa\,ri}(\nu,\bar\nu,\omega)
    \right]
    Q^+_{ia}(\bar\nu,\omega)
    \label{eq:U_eh_second_def}
\end{equation}
which gives
\begin{equation}
    \begin{split}
    U^{\eh,+,(2)}_{pr,ia}(\nu,\omega)
    & =
    \sum_{jb}
    \frac{\mel{pb}{}{ij}\mel{aj}{}{rb}}
    {\nu-(\eps_i+\eps_j-\eps_b)-\ii0^+}
    \\
    &-
    \sum_{jb}
    \frac{\mel{pj}{}{ib}\mel{ab}{}{rj}}
    {\nu+\omega-(\eps_a+\eps_b-\eps_j)+\ii0^+}
    \\
    &-
    \sum_{j<k}
    \frac{\mel{pa}{}{jk}\mel{jk}{}{ri}}
    {\nu-(\eps_j+\eps_k-\eps_a)-\ii0^+}
    \\
    &+
    \sum_{b<c}
    \frac{\mel{pa}{}{bc}\mel{bc}{}{ri}}
    {\nu+\omega-(\eps_b+\eps_c-\eps_i)+\ii0^+}
    \end{split}
    \label{eq:U_eh_second_explicit}
\end{equation}
The corresponding right correction is
\begin{equation}
    V^{\eh,+,(2)}_{ia,sq}(\nu',\omega)
    =
    \frac{1}{2\pi}
    \int\dd{\bar\nu}
    Q^+_{ia}(\bar\nu,\omega)
    \left[
    \ii\Gamma^{\eh,(2)}_{iq\,as}(\bar\nu,\nu',\omega)
    \right]
    \label{eq:V_eh_second_def}
\end{equation}
with
\begin{equation}
    \begin{split}
    V^{\eh,+,(2)}_{ia,sq}(\nu',\omega)
    & =
    \sum_{kc}
    \frac{\mel{sc}{}{ik}\mel{ak}{}{qc}}
    {\nu'-(\eps_i+\eps_k-\eps_c)-\ii0^+}
    \\
    &-
    \sum_{kc}
    \frac{\mel{sk}{}{ic}\mel{ac}{}{qk}}
    {\nu'+\omega-(\eps_a+\eps_c-\eps_k)+\ii0^+}
    \\
    &-
    \sum_{k<l}
    \frac{\mel{sa}{}{kl}\mel{kl}{}{qi}}
    {\nu'-(\eps_k+\eps_l-\eps_a)-\ii0^+}
    \\
    &+
    \sum_{c<d}
    \frac{\mel{sa}{}{cd}\mel{cd}{}{qi}}
    {\nu'+\omega-(\eps_c+\eps_d-\eps_i)+\ii0^+}
    \end{split}
    \label{eq:V_eh_second_explicit}
\end{equation}

For the $\text{vo}$ sector,
\begin{equation}
    U^{\eh,-,(2)}_{pr,ia}(\nu,\omega)
    =
    \frac{1}{2\pi}
    \int\dd{\bar\nu}
    \left[
    \ii\Gamma^{\eh,(2)}_{pi\,ra}(\nu,\bar\nu,\omega)
    \right]
    Q^-_{ia}(\bar\nu,\omega)
    \label{eq:U_eh_second_minus_def}
\end{equation}
and the contour integrations give
\begin{equation}
    \begin{split}
    U^{\eh,-,(2)}_{pr,ia}(\nu,\omega)
    & =
    \sum_{kc}
    \frac{\mel{pc}{}{ak}\mel{ik}{}{rc}}
    {\nu+\omega-(\eps_i+\eps_k-\eps_c)-\ii0^+}
    \\
    &-
    \sum_{kc}
    \frac{\mel{pk}{}{ac}\mel{ic}{}{rk}}
    {\nu-(\eps_a+\eps_c-\eps_k)+\ii0^+}
    \\
    &+
    \sum_{c<d}
    \frac{\mel{pi}{}{cd}\mel{cd}{}{ra}}
    {\nu-(\eps_c+\eps_d-\eps_i)+\ii0^+}
    \\
    &-
    \sum_{k<l}
    \frac{\mel{pi}{}{kl}\mel{kl}{}{ra}}
    {\nu+\omega-(\eps_k+\eps_l-\eps_a)-\ii0^+}
    \end{split}
    \label{eq:U_eh_second_minus_explicit}
\end{equation}
Likewise,
\begin{equation}
    V^{\eh,-,(2)}_{ia,sq}(\nu',\omega)
    =
    \frac{1}{2\pi}
    \int\dd{\bar\nu}
    Q^-_{ia}(\bar\nu,\omega)
    \left[
    \ii\Gamma^{\eh,(2)}_{aq\,is}(\bar\nu,\nu',\omega)
    \right]
    \label{eq:V_eh_second_minus_def}
\end{equation}
with
\begin{equation}
    \begin{split}
    V^{\eh,-,(2)}_{ia,sq}(\nu',\omega)
    & =
    \sum_{kc}
    \frac{\mel{ak}{}{sc}\mel{qc}{}{ik}}
    {\nu'+\omega-(\eps_i+\eps_k-\eps_c)-\ii0^+}
    \\
    &-
    \sum_{kc}
    \frac{\mel{ac}{}{sk}\mel{qk}{}{ic}}
    {\nu'-(\eps_a+\eps_c-\eps_k)+\ii0^+}
    \\
    &-
    \sum_{k<l}
    \frac{\mel{aq}{}{kl}\mel{kl}{}{is}}
    {\nu'+\omega-(\eps_k+\eps_l-\eps_a)-\ii0^+}
    \\
    &+
    \sum_{c<d}
    \frac{\mel{aq}{}{cd}\mel{cd}{}{is}}
    {\nu'-(\eps_c+\eps_d-\eps_i)+\ii0^+}
    \end{split}
    \label{eq:V_eh_second_minus_explicit}
\end{equation}
The $+$ and $-$ components contain $\hhp$ and $\pph$ poles,
as anticipated for the second-order coupling corrections.

%%%%%%%%%%%%%%%%%%%%%%%%%
\subsection{$\text{oo}$ and $\text{vv}$ sectors}
%%%%%%%%%%%%%%%%%%%%%%%%%

The $\text{vv}$ sector selects the poles of $\Gamma^{\eh,(2)}$ lying in the
half-plane opposite to the two virtual one-particle poles.
For the first ordering in Eq.~\eqref{eq:Phi_eh_third_convolution},
\begin{widetext}
\begin{equation}
    \begin{split}
    \left\{
    \left[\ii\Gamma^{\eh,(2)}\right]
    \circ\mathcal L_0^{\eh,\text{vv}}
    \circ\left[\ii\Gamma^{\eh,(1)}\right]
    \right\}_{pqrs}
    & =
    \sum_{abkc}
    \mel{pk}{}{bc}\mel{ac}{}{rk}\mel{qb}{}{sa}
    \frac{1}{\nu+\omega-(\eps_a+\eps_c-\eps_k)+\ii0^+}
    \frac{1}{\nu-(\eps_b+\eps_c-\eps_k)+\ii0^+}
    \\
    &-
    \sum_{ab}\sum_{k<l}
    \mel{pa}{}{kl}\mel{kl}{}{rb}\mel{qb}{}{sa}
    \frac{1}{\nu-(\eps_k+\eps_l-\eps_a)-\ii0^+}
    \frac{1}{\nu+\omega-(\eps_k+\eps_l-\eps_b)-\ii0^+}
    \end{split}
    \label{eq:Phi_eh_third_vv_left}
\end{equation}
and the opposite ordering gives
\begin{equation}
    \begin{split}
    \left\{
    \left[\ii\Gamma^{\eh,(1)}\right]
    \circ\mathcal L_0^{\eh,\text{vv}}
    \circ\left[\ii\Gamma^{\eh,(2)}\right]
    \right\}_{pqrs}
    & =
    \sum_{abkc}
    \mel{pa}{}{rb}\mel{bc}{}{sk}\mel{qk}{}{ac}
    \frac{1}{\nu'+\omega-(\eps_a+\eps_c-\eps_k)+\ii0^+}
    \frac{1}{\nu'-(\eps_b+\eps_c-\eps_k)+\ii0^+}
    \\
    &-
    \sum_{ab}\sum_{k<l}
    \mel{pa}{}{rb}\mel{bq}{}{kl}\mel{kl}{}{as}
    \frac{1}{\nu'-(\eps_k+\eps_l-\eps_a)-\ii0^+}
    \frac{1}{\nu'+\omega-(\eps_k+\eps_l-\eps_b)-\ii0^+}
    \end{split}
    \label{eq:Phi_eh_third_vv_right}
\end{equation}

The $\text{oo}$ sector selects the opposite causal poles of the second-order kernel.
For the first ordering,
\begin{equation}
    \begin{split}
    \left\{
    \left[\ii\Gamma^{\eh,(2)}\right]
    \circ\mathcal L_0^{\eh,\text{oo}}
    \circ\left[\ii\Gamma^{\eh,(1)}\right]
    \right\}_{pqrs}
    & =
    \sum_{ijkc}
    \mel{pc}{}{jk}\mel{ik}{}{rc}\mel{qj}{}{si}
    \frac{1}{\nu+\omega-(\eps_i+\eps_k-\eps_c)-\ii0^+}
    \frac{1}{\nu-(\eps_j+\eps_k-\eps_c)-\ii0^+}
    \\
    &-
    \sum_{ij}\sum_{c<d}
    \mel{pi}{}{cd}\mel{cd}{}{rj}\mel{qj}{}{si}
    \frac{1}{\nu-(\eps_c+\eps_d-\eps_i)+\ii0^+}
    \frac{1}{\nu+\omega-(\eps_c+\eps_d-\eps_j)+\ii0^+}
    \end{split}
    \label{eq:Phi_eh_third_oo_left}
\end{equation}
and for the opposite ordering,
\begin{equation}
    \begin{split}
    \left\{
    \left[\ii\Gamma^{\eh,(1)}\right]
    \circ\mathcal L_0^{\eh,\text{oo}}
    \circ\left[\ii\Gamma^{\eh,(2)}\right]
    \right\}_{pqrs}
    & =
    \sum_{ijkc}
    \mel{pi}{}{rj}\mel{jk}{}{sc}\mel{qc}{}{ik}
    \frac{1}{\nu'+\omega-(\eps_i+\eps_k-\eps_c)-\ii0^+}
    \frac{1}{\nu'-(\eps_j+\eps_k-\eps_c)-\ii0^+}
    \\
    &-
    \sum_{ij}\sum_{c<d}
    \mel{pi}{}{rj}\mel{jq}{}{cd}\mel{cd}{}{is}
    \frac{1}{\nu'-(\eps_c+\eps_d-\eps_i)+\ii0^+}
    \frac{1}{\nu'+\omega-(\eps_c+\eps_d-\eps_j)+\ii0^+}
    \end{split}
    \label{eq:Phi_eh_third_oo_right}
\end{equation}
\end{widetext}
These expressions are deliberately left as products of charged-state
resolvents.
No partial-fraction decomposition is introduced, so no denominator involving
$\eps_a-\eps_b$ or $\eps_i-\eps_j$ appears.

%%%%%%%%%%%%%%%%%%%%%%%%%
\subsection{Three-interaction term}
%%%%%%%%%%%%%%%%%%%%%%%%%

Because all three vertices in the last term of
Eq.~\eqref{eq:Phi_eh_third_convolution} are instantaneous, only the $\text{ov}$ and
$\text{vo}$ sectors survive the internal frequency integrations.
The first-order neutral interaction matrix is
\begin{equation}
    M^{\eh,(1)}
    =
    \begin{pmatrix}
    A^{(1)}&B^{(1)}
    \\
    B^{(1)}&A^{(1)}
    \end{pmatrix}
    \label{eq:M_eh_first}
\end{equation}
with
\begin{align}
    A^{(1)}_{ia,jb}
    &=
    \mel{aj}{}{ib}
    \\
    B^{(1)}_{ia,jb}
    &=
    \mel{ab}{}{ij}
\end{align}
and
\begin{equation}
    R^{\eh,(1)}
    =
    R^{\eh,(0)}M^{\eh,(1)}R^{\eh,(0)}
    \label{eq:R_eh_first}
\end{equation}
Thus, the complete contribution from the three instantaneous interactions is
\begin{equation}
    U^{\eh,(1)}R^{\eh,(0)}M^{\eh,(1)}R^{\eh,(0)}V^{\eh,(1)}
    \label{eq:Phi_eh_third_three_bare}
\end{equation}
For example, its forward-forward block is
\begin{equation}
    \begin{split}
    &\left[
    U^{\eh,(1)}R^{\eh,(0)}M^{\eh,(1)}R^{\eh,(0)}V^{\eh,(1)}
    \right]^{++}_{pqrs}
    \\
    &\quad=
    \sum_{iajb}
    \frac{
    \mel{pa}{}{ri}\mel{aj}{}{ib}\mel{qj}{}{sb}
    }
    {[\omega-(\eps_a-\eps_i)+\ii0^+]
     [\omega-(\eps_b-\eps_j)+\ii0^+]}
    \end{split}
    \label{eq:Phi_eh_third_three_bare_pp}
\end{equation}

%%%%%%%%%%%%%%%%%%%%%%%%%
\subsection{Complete direct eh result through third order}
%%%%%%%%%%%%%%%%%%%%%%%%%

The complete result is most transparently written by retaining the four
orbital sectors of Eq.~\eqref{eq:L0_sector_decomposition},
\begin{equation}
    \begin{split}
    \ii\Phi^{\eh,(3)}
    &=
    \sum_{x\in\{\text{oo},\text{ov},\text{vo},\text{vv}\}}
    \Big(
    \left[\ii\Gamma^{\eh,(2)}\right]
    \circ\mathcal L_0^{\eh,x}
    \circ\left[\ii\Gamma^{\eh,(1)}\right]
    \\
    &\qquad+
    \left[\ii\Gamma^{\eh,(1)}\right]
    \circ\mathcal L_0^{\eh,x}
    \circ\left[\ii\Gamma^{\eh,(2)}\right]
    \Big)
    \\
    &+
    U^{\eh,(1)}R^{\eh,(0)}M^{\eh,(1)}R^{\eh,(0)}V^{\eh,(1)}
    \end{split}
    \label{eq:Phi_eh_third_final}
\end{equation}
The $\text{ov}$ and $\text{vo}$ subset of the first two lines is equivalently
\begin{equation}
    U^{\eh,(2)}R^{\eh,(0)}V^{\eh,(1)}
    +
    U^{\eh,(1)}R^{\eh,(0)}V^{\eh,(2)}
    \label{eq:Phi_eh_third_ovvo_resolvent}
\end{equation}
while the $\text{oo}$ and $\text{vv}$ terms are given explicitly in
Eqs.~\eqref{eq:Phi_eh_third_vv_left}--\eqref{eq:Phi_eh_third_oo_right}.
The three-interaction term is Eq.~\eqref{eq:Phi_eh_third_three_bare} and is
equivalent to $U^{\eh,(1)}R^{\eh,(1)}V^{\eh,(1)}$.
No additional orbital sector or third-order topology occurs in the direct
$\eh$ reducible vertex.

\section{pp vertex through third order}
\label{sec:pp_third}
%%%%%%%%%%%%%%%%%%%%%%%%%

The $\pp$ construction uses the unique-pair convention of
Sec.~\ref{sec:frequencies} and is treated in the native variables
\begin{equation}
    x=\nu_\pp
    \qquad
    y=\nu'_\pp
    \qquad
    \Omega=\Omega_\pp
    \label{eq:pp_native_variables_third}
\end{equation}
for which the four external fermionic lines carry
$\Omega+x$, $-x$, $\Omega+y$, and $-y$.
In terms of the direct-$\eh$ variables,
\begin{equation}
    \nu=\Omega+x
    \qquad
    \nu'=-y
    \qquad
    \omega=y-x
    \label{eq:pp_inverse_frequency_map_third}
\end{equation}
and therefore $\nu'-\nu=-(\Omega+x+y)$.

%%%%%%%%%%%%%%%%%%%%%%%%%
\subsection{Perturbative BSE expansion}
%%%%%%%%%%%%%%%%%%%%%%%%%

In each internal $\pp$ contraction, $\circ$ includes the average over the
two orbital and relative-frequency orientations of each unique pair.
Equivalently, an ordered-orientation sum carries its usual symmetry weight
$1/2$ for each internal pair propagation.
After the orientation average, matrix contractions use $a<b$ and $i<j$,
and the mixed pair $ia$ is counted once.
The oriented propagator defined below does not itself contain this weight.

At third order, the $\pp$ reducible vertex contains the same three
perturbative contributions as the direct-$\eh$ vertex.
Suppressing the explicit internal-frequency integrations with the convolution
symbol $\circ$, and using the reduced native pair propagator
$\mathcal L_0^\pp$ defined below,
\begin{equation}
    \begin{split}
    \ii\Phi^{\pp,(3)}
    & =
    \left[\ii\Gamma^{\pp,(2)}\right]
    \circ\mathcal L_0^\pp
    \circ\left[\ii\Gamma^{\pp,(1)}\right]
    \\
    &+
    \left[\ii\Gamma^{\pp,(1)}\right]
    \circ\mathcal L_0^\pp
    \circ\left[\ii\Gamma^{\pp,(2)}\right]
    \\
    &+
    \left[\ii\Gamma^{\pp,(1)}\right]
    \circ\mathcal L_0^\pp
    \circ\left[\ii\Gamma^{\pp,(1)}\right]
    \circ\mathcal L_0^\pp
    \circ\left[\ii\Gamma^{\pp,(1)}\right]
    \end{split}
    \label{eq:Phi_pp_third_convolution}
\end{equation}
The first two terms contain the frequency-dependent second-order irreducible
kernel on opposite sides of the internal propagation.
The last term contains three instantaneous first-order interactions.

%%%%%%%%%%%%%%%%%%%%%%%%%
\subsection{Native kernel and orbital decomposition}
%%%%%%%%%%%%%%%%%%%%%%%%%

Using
$\Gamma^{\pp,(2)}=\Phi^{\eh,(2)}+\Phi^{\teh,(2)}$
and the map above gives
\begin{equation}
\begin{split}
    \ii\Gamma^{\pp,(2)}_{pqrs}(x,y,\Omega)
    &=
    \sum_{kc}
    \frac{
    \mel{pc}{}{rk}\mel{qk}{}{sc}
    }
    {y-x-(\eps_c-\eps_k)+\ii0^+}
    \\
    & -
    \sum_{kc}
    \frac{
    \mel{pk}{}{rc}\mel{qc}{}{sk}
    }
    {y-x+(\eps_c-\eps_k)-\ii0^+}
    \\
    &+
    \sum_{kc}
    \frac{
    \mel{pc}{}{sk}\mel{qk}{}{rc}
    }
    {\Omega+x+y+(\eps_c-\eps_k)-\ii0^+}
    \\
    & -
    \sum_{kc}
    \frac{
    \mel{pk}{}{sc}\mel{qc}{}{rk}
    }
    {\Omega+x+y-(\eps_c-\eps_k)+\ii0^+}
\end{split}
\label{eq:Gamma_pp_second_native_third}
\end{equation}

It is convenient to define the reduced
native pair propagator
\begin{equation}
    \mathcal L_{0,tu}^{\pp}(y,\Omega)
    =
    \ii G_{0,t}(\Omega+y)G_{0,u}(-y)
    \label{eq:L0_pp_reduced_oriented}
\end{equation}
and decompose
\begin{equation}
    \mathcal L_0^\pp
    =
    \mathcal L_0^{\pp,\text{oo}}
    +
    \mathcal L_0^{\pp,\text{ov}}
    +
    \mathcal L_0^{\pp,\text{vo}}
    +
    \mathcal L_0^{\pp,\text{vv}}
    \label{eq:L0_pp_sector_decomposition}
\end{equation}
For an instantaneous kernel,
\begin{align}
    \frac{1}{2\pi}\int\dd{y}\,
    \mathcal L_{0,ab}^{\pp}(y,\Omega)
    &=
    \frac{1}{\Omega-(\eps_a+\eps_b)+\ii0^+}
    \label{eq:pp_vv_integral}
    \\
    \frac{1}{2\pi}\int\dd{y}\,
    \mathcal L_{0,ij}^{\pp}(y,\Omega)
    &=
    -\frac{1}{\Omega-(\eps_i+\eps_j)-\ii0^+}
    \label{eq:pp_oo_integral}
    \\
    \frac{1}{2\pi}\int\dd{y}\,
    \mathcal L_{0,ia}^{\pp}(y,\Omega)
    &=0
    \\
    \frac{1}{2\pi}\int\dd{y}\,
    \mathcal L_{0,ai}^{\pp}(y,\Omega)
    &=0
    \label{eq:pp_mixed_integrals_zero}
\end{align}
Thus, the $\text{oo}$ and $\text{vv}$ sectors generate the ordinary pair resolvent, whereas
the isolated $\text{ov}$ and $\text{vo}$ sectors vanish.
A dynamic $\Gamma^{\pp,(2)}$ supplies an additional pole and activates the
mixed sectors at third order.

%%%%%%%%%%%%%%%%%%%%%%%%%
\subsection{$\text{oo}$ and $\text{vv}$ sectors}
%%%%%%%%%%%%%%%%%%%%%%%%%

The $\text{oo}$ and $\text{vv}$ parts of the first two terms in
Eq.~\eqref{eq:Phi_pp_third_convolution} can be reorganized as corrections to
the coupling vertices.
For a unique virtual pair $a<b$, define the two oriented normalized factors
\begin{align}
    Q_{ab}^{\pp,+}(y,\Omega)
    &=
    \ii\left[G_{0,a}(\Omega+y)+G_{0,b}(-y)\right]
    \\
    Q_{ba}^{\pp,+}(y,\Omega)
    &=
    \ii\left[G_{0,b}(\Omega+y)+G_{0,a}(-y)\right]
    \label{eq:Q_pp_plus_orientations}
\end{align}
and, for a unique occupied pair $i<j$,
\begin{align}
    Q_{ij}^{\pp,-}(y,\Omega)
    &=
    -\ii\left[G_{0,i}(\Omega+y)+G_{0,j}(-y)\right]
    \\
    Q_{ji}^{\pp,-}(y,\Omega)
    &=
    -\ii\left[G_{0,j}(\Omega+y)+G_{0,i}(-y)\right]
    \label{eq:Q_pp_minus_orientations}
\end{align}
Each oriented factor integrates to unity.
The antisymmetric unique-pair projection must be applied to the full
kernel--propagator product.  For example,
\begin{equation}
\begin{split}
    U_{pq,ab}^{\pp,(2)}(x,\Omega)
    &=
    \frac12\int\frac{\dd{y}}{2\pi}
    \Big[
    \ii\Gamma_{pqab}^{\pp,(2)}(x,y,\Omega)
    Q_{ab}^{\pp,+}(y,\Omega)
    \\
    &\qquad-
    \ii\Gamma_{pqba}^{\pp,(2)}(x,y,\Omega)
    Q_{ba}^{\pp,+}(y,\Omega)
    \Big]
\end{split}
\label{eq:U_pp_second_projection_ab}
\end{equation}
with the analogous expression for $ij$ and for the right coupling vertex.
At first order, the same projection gives
\begin{align}
	U_{pq,ab}^{\pp,(1)} & = \mel{pq}{}{ab} 
	\\
	U_{pq,ij}^{\pp,(1)} & = \mel{pq}{}{ij}
\end{align}
fixing the normalization.

The explicit left coupling vertices are
\begin{equation}
\begin{split}
U_{pq,ab}^{\pp,(2)}(x,\Omega)
&=
\sum_{kc}
\frac{\mel{pc}{}{bk}\mel{qk}{}{ac}}
{x+(\eps_a+\eps_c-\eps_k)-\ii0^+}
\\
&-
\sum_{kc}
\frac{\mel{pc}{}{ak}\mel{qk}{}{bc}}
{x+(\eps_b+\eps_c-\eps_k)-\ii0^+}
\\
&+
\sum_{kc}
\frac{\mel{pk}{}{ac}\mel{qc}{}{bk}}
{\Omega+x-(\eps_a+\eps_c-\eps_k)+\ii0^+}
\\
&-
\sum_{kc}
\frac{\mel{pk}{}{bc}\mel{qc}{}{ak}}
{\Omega+x-(\eps_b+\eps_c-\eps_k)+\ii0^+}
\end{split}
\label{eq:U_pp_second_ab}
\end{equation}
and
\begin{equation}
\begin{split}
U_{pq,ij}^{\pp,(2)}(x,\Omega)
&=
-
\sum_{kc}
\frac{\mel{pc}{}{ik}\mel{qk}{}{jc}}
{\Omega+x-(\eps_i+\eps_k-\eps_c)-\ii0^+}
\\
&+
\sum_{kc}
\frac{\mel{pc}{}{jk}\mel{qk}{}{ic}}
{\Omega+x-(\eps_j+\eps_k-\eps_c)-\ii0^+}
\\
&-
\sum_{kc}
\frac{\mel{pk}{}{jc}\mel{qc}{}{ik}}
{x+(\eps_i+\eps_k-\eps_c)+\ii0^+}
\\
&+
\sum_{kc}
\frac{\mel{pk}{}{ic}\mel{qc}{}{jk}}
{x+(\eps_j+\eps_k-\eps_c)+\ii0^+}
\end{split}
\label{eq:U_pp_second_ij}
\end{equation}
They satisfy
$U_{pq,ba}^{\pp,(2)}=-U_{pq,ab}^{\pp,(2)}$ and
$U_{pq,ji}^{\pp,(2)}=-U_{pq,ij}^{\pp,(2)}$.

The corresponding right coupling vertices are
\begin{equation}
\begin{split}
V_{ab,rs}^{\pp,(2)}(y,\Omega)
&=
\sum_{kc}
\frac{\mel{ac}{}{rk}\mel{bk}{}{sc}}
{\Omega+y-(\eps_a+\eps_c-\eps_k)+\ii0^+}
\\
& -
\sum_{kc}
\frac{\mel{bc}{}{rk}\mel{ak}{}{sc}}
{\Omega+y-(\eps_b+\eps_c-\eps_k)+\ii0^+}
\\
&+
\sum_{kc}
\frac{\mel{ac}{}{sk}\mel{bk}{}{rc}}
{y+(\eps_a+\eps_c-\eps_k)-\ii0^+}
\\
& -
\sum_{kc}
\frac{\mel{bc}{}{sk}\mel{ak}{}{rc}}
{y+(\eps_b+\eps_c-\eps_k)-\ii0^+}
\end{split}
\label{eq:V_pp_second_ab}
\end{equation}
and
\begin{equation}
\begin{split}
V_{ij,rs}^{\pp,(2)}(y,\Omega)
&=
-
\sum_{kc}
\frac{\mel{ik}{}{sc}\mel{jc}{}{rk}}
{y+(\eps_i+\eps_k-\eps_c)+\ii0^+}
\\
&+
\sum_{kc}
\frac{\mel{ic}{}{rk}\mel{jk}{}{sc}}
{y+(\eps_j+\eps_k-\eps_c)+\ii0^+}
\\
&-
\sum_{kc}
\frac{\mel{ik}{}{rc}\mel{jc}{}{sk}}
{\Omega+y-(\eps_i+\eps_k-\eps_c)-\ii0^+}
\\
&+
\sum_{kc}
\frac{\mel{ic}{}{sk}\mel{jk}{}{rc}}
{\Omega+y-(\eps_j+\eps_k-\eps_c)-\ii0^+}
\end{split}
\label{eq:V_pp_second_ij}
\end{equation}
with
$V_{ba,rs}^{\pp,(2)}=-V_{ab,rs}^{\pp,(2)}$ and
$V_{ji,rs}^{\pp,(2)}=-V_{ij,rs}^{\pp,(2)}$.
All nontrivial poles are again of $\pph$ or $\hhp$ type.

%%%%%%%%%%%%%%%%%%%%%%%%%
\subsection{$\text{ov}$ and $\text{vo}$ sectors}
%%%%%%%%%%%%%%%%%%%%%%%%%

The $\text{ov}$ and $\text{vo}$ parts of the first two terms in
Eq.~\eqref{eq:Phi_pp_third_convolution} do not belong to the ordinary pair
resolvent and are evaluated explicitly.
For the $\text{ov}$ orientation $ia$, both one-particle poles lie in the upper
half-plane, so the first and fourth terms of
Eq.~\eqref{eq:Gamma_pp_second_native_third} contribute.
For the $\text{vo}$ orientation $ai$, both lie in the lower half-plane, so the
second and third terms contribute.
Using $\mel{ai}{}{rs}=-\mel{ia}{}{rs}$, the two orientations give equal
contributions and cancel the factor $1/2$ associated with summing both pair orientations.
The resulting first ordering is
\begin{widetext}
\begin{equation}
\begin{split}
\left\{
\left[\ii\Gamma^{\pp,(2)}\right]
\circ\left(\mathcal L_0^{\pp,\text{ov}}+\mathcal L_0^{\pp,\text{vo}}\right)
\circ\left[\ii\Gamma^{\pp,(1)}\right]
\right\}_{pqrs}
&=
-
\sum_{iakc}
\frac{
\mel{pc}{}{ik}\mel{qk}{}{ac}\mel{ia}{}{rs}
}
{[\Omega+x-(\eps_i+\eps_k-\eps_c)-\ii0^+]
 [x+(\eps_a+\eps_c-\eps_k)-\ii0^+]}
\\
&+
\sum_{iakc}
\frac{
\mel{pk}{}{ac}\mel{qc}{}{ik}\mel{ia}{}{rs}
}
{[x+(\eps_i+\eps_k-\eps_c)+\ii0^+]
 [\Omega+x-(\eps_a+\eps_c-\eps_k)+\ii0^+]}
\end{split}
\label{eq:Phi_pp_third_mixed_left}
\end{equation}
and the reversed ordering is
\begin{equation}
\begin{split}
\left\{
\left[\ii\Gamma^{\pp,(1)}\right]
\circ\left(\mathcal L_0^{\pp,\text{ov}}+\mathcal L_0^{\pp,\text{vo}}\right)
\circ\left[\ii\Gamma^{\pp,(2)}\right]
\right\}_{pqrs}
&=
-
\sum_{iakc}
\frac{
\mel{pq}{}{ia}\mel{ik}{}{rc}\mel{ac}{}{sk}
}
{[\Omega+y-(\eps_i+\eps_k-\eps_c)-\ii0^+]
 [y+(\eps_a+\eps_c-\eps_k)-\ii0^+]}
\\
&+
\sum_{iakc}
\frac{
\mel{pq}{}{ia}\mel{ik}{}{sc}\mel{ac}{}{rk}
}
{[y+(\eps_i+\eps_k-\eps_c)+\ii0^+]
 [\Omega+y-(\eps_a+\eps_c-\eps_k)+\ii0^+]}
\end{split}
\label{eq:Phi_pp_third_mixed_right}
\end{equation}
\end{widetext}
No partial fractioning is required, and only charged-state energy combinations
appear in these denominators.

%%%%%%%%%%%%%%%%%%%%%%%%%
\subsection{Three-interaction term}
%%%%%%%%%%%%%%%%%%%%%%%%%

Because all three vertices in the last term of
Eq.~\eqref{eq:Phi_pp_third_convolution} are instantaneous, only the
$\text{oo}$ and $\text{vv}$ sectors survive the internal frequency integrations.
The first-order pair interaction matrix is
\begin{equation}
    M^{\pp,(1)}
    =
    \begin{pmatrix}
    \mel{ab}{}{cd}&\mel{ab}{}{ij}
    \\
    \mel{ij}{}{cd}&\mel{ij}{}{kl}
    \end{pmatrix}
    \label{eq:M_pp_first}
\end{equation}
with
\begin{equation}
    R^{\pp,(1)}
    =
    R^{\pp,(0)}M^{\pp,(1)}R^{\pp,(0)}
    \label{eq:R_pp_first}
\end{equation}
Thus, the complete contribution from the three instantaneous interactions is
\begin{equation}
    U^{\pp,(1)}R^{\pp,(0)}M^{\pp,(1)}R^{\pp,(0)}V^{\pp,(1)}
    \label{eq:Phi_pp_third_three_bare}
\end{equation}
For example, the pure $\text{vv}$ contribution is
\begin{equation}
\begin{split}
    &\left[
    U^{\pp,(1)}R^{\pp,(0)}M^{\pp,(1)}R^{\pp,(0)}V^{\pp,(1)}
    \right]_{pqrs}^{\text{vv}}
    \\
    &\quad=
    \sum_{\substack{a<b\\c<d}}
    \frac{
    \mel{pq}{}{ab}\mel{ab}{}{cd}\mel{cd}{}{rs}
    }
    {[\Omega-(\eps_a+\eps_b)+\ii0^+]
     [\Omega-(\eps_c+\eps_d)+\ii0^+]}
\end{split}
\label{eq:Phi_pp_third_central_vv}
\end{equation}
The $\text{vv}\leftrightarrow\text{oo}$ and $\text{oo}\to\text{oo}$ blocks follow from the remaining
blocks of $M^{\pp,(1)}$.

%%%%%%%%%%%%%%%%%%%%%%%%%
\subsection{Complete pp result through third order}
%%%%%%%%%%%%%%%%%%%%%%%%%

The complete result is most transparently written by retaining the four
orbital sectors of Eq.~\eqref{eq:L0_pp_sector_decomposition},
\begin{equation}
    \begin{split}
    \ii\Phi^{\pp,(3)}
    &=
    \sum_{x\in\{\text{oo},\text{ov},\text{vo},\text{vv}\}}
    \Big(
    \left[\ii\Gamma^{\pp,(2)}\right]
    \circ\mathcal L_0^{\pp,x}
    \circ\left[\ii\Gamma^{\pp,(1)}\right]
    \\
    &\qquad+
    \left[\ii\Gamma^{\pp,(1)}\right]
    \circ\mathcal L_0^{\pp,x}
    \circ\left[\ii\Gamma^{\pp,(2)}\right]
    \Big)
    \\
    &+
    U^{\pp,(1)}R^{\pp,(0)}M^{\pp,(1)}R^{\pp,(0)}V^{\pp,(1)}
    \end{split}
    \label{eq:Phi_pp_third_final}
\end{equation}
The $\text{oo}$ and $\text{vv}$ subset of the first two lines is equivalently
\begin{equation}
    U^{\pp,(2)}R^{\pp,(0)}V^{\pp,(1)}
    +
    U^{\pp,(1)}R^{\pp,(0)}V^{\pp,(2)}
    \label{eq:Phi_pp_third_oovv_resolvent}
\end{equation}
while the $\text{ov}$ and $\text{vo}$ terms are given explicitly in
Eqs.~\eqref{eq:Phi_pp_third_mixed_left} and
\eqref{eq:Phi_pp_third_mixed_right}.
The three-interaction term is Eq.~\eqref{eq:Phi_pp_third_three_bare} and is
equivalent to $U^{\pp,(1)}R^{\pp,(1)}V^{\pp,(1)}$.
No additional orbital sector or third-order topology occurs in the primitive
$\pp$ reducible vertex.


\bibliography{multiresolvent}

\end{document}

